\documentclass[10pt,aspectratio=169,professionalfonts]{beamer}

% ===== 编译建议 =====
% 推荐使用 XeLaTeX 编译

% ===== 主题与颜色 =====
\usetheme{Madrid}
\usecolortheme{default}

% ===== 中文支持 =====
\usepackage[UTF8]{ctex}

% ===== 常用数学与图形 =====
\usepackage{amsmath,amssymb,bm}
\usepackage{booktabs}
\usepackage{graphicx}
\usepackage{hyperref}
\usepackage{ragged2e}

% ===== 页面信息（按需修改） =====
\title{ME-gPC and WENO}
\author{Li Yunfan}
\institute{Zhejiang University}
\date{2026.3.11}

% ===== 每节开头自动目录 =====
\AtBeginSection[]
{
	\begin{frame}{Contents}
		\tableofcontents[currentsection]
	\end{frame}
}

\begin{document}

% ===== 封面 =====
\begin{frame}
	\titlepage
\end{frame}

% ===== 总目录 =====
\begin{frame}{Contents}
	\tableofcontents
\end{frame}

% ==========================
\section{Multi-element Generalized Polynomial Chaos}

\begin{frame}{Introduction}
    \justifying
	Multi-element Generalized Polynomial Chaos is a method for 
    uncertainty quantification that divides the input space into 
    multiple elements and applies polynomial chaos expansions within 
    each element. This approach allows for better handling of non-smooth 
    functions and localized features in the solution space.
    The key idea is to partition the random input space 
    into smaller subdomains.
\end{frame}

\begin{frame}{Decomposition of Random Space}
    \justifying
	To describe a problem involving uncertainty, 
    we first need to specify the 
    underlying probability space \((\Omega,\mathcal{F},P)\).
    However, the probability space is often absatract and 
    not directly usable for numerical methods, 
	which causes difficulties in the decomposition.\cite{WAN2005617}\cite{Wan2006}\cite{FOO20089572}
    
    
\end{frame}

% ==========================
\section{方法介绍}

\begin{frame}{核心思路}
	\begin{enumerate}
		\item 输入与建模假设
		\item 算法流程与关键步骤
		\item 复杂度或理论性质
	\end{enumerate}
\end{frame}

\begin{frame}{关键公式}
	示例公式：
	\[
		\hat{\theta} = \arg\min_{\theta} \; \mathcal{L}(\theta)
	\]

	\vspace{0.6em}
	其中 \(\mathcal{L}(\theta)\) 表示目标函数。
\end{frame}

\begin{frame}{图表示例}
	\begin{figure}
		\centering
		% 将 example-image 替换为你的图片路径
		\includegraphics[width=0.72\linewidth]{example-image}
		\caption{示意图}
	\end{figure}
\end{frame}

% ==========================
\section{实验结果}

\begin{frame}{实验设置}
	\begin{itemize}
		\item 数据集：XXX
		\item 评价指标：XXX
		\item 对比方法：XXX
	\end{itemize}
\end{frame}

\begin{frame}{结果对比}
	\begin{table}
		\centering
		\begin{tabular}{lccc}
			\toprule
			方法 & 指标1 & 指标2 & 指标3 \\
			\midrule
			Baseline A & 0.81 & 0.75 & 0.68 \\
			Baseline B & 0.84 & 0.77 & 0.71 \\
			Ours & \textbf{0.89} & \textbf{0.82} & \textbf{0.76} \\
			\bottomrule
		\end{tabular}
		\caption{实验结果示例}
	\end{table}
\end{frame}

% ==========================
\section{总结与展望}

\begin{frame}{总结}
	\begin{itemize}
		\item 工作内容与主要结论
		\item 方法优势与适用范围
		\item 局限性与未来工作方向
	\end{itemize}
\end{frame}

\begin{frame}{致谢}
	\centering
	\Large 谢谢！\\
	\vspace{0.8em}
	\normalsize 欢迎提问。
\end{frame}

% ===== 参考文献 =====
\begin{frame}[allowframebreaks]{References}
	\bibliographystyle{plain}
	\bibliography{refs}
\end{frame}

\end{document}
